\section{Eliminating the payment rule}
\label{sec:elimination}

Feasibility in Section~\ref{sec:model} couples two unknowns, the record $\enc$
and the decoders $\pay_i$. Products of the two make the natural formulation
bilinear. Our first result removes $\pay_i$ entirely.

\begin{theorem}[Payment-decoder elimination]
\label{thm:elim}
Fix a record $\enc$ and an agent $i$. There exists a bounded decoder
$\pay_i \in [0,1]^K$ satisfying \eqref{eq:margin} if and only if
\[
  \min_{\lambda \in \simplex(\{0,1\}^{n-1})}
  \tfrac{1}{2}\bigl\|\,\E_{\lambda}\bigl[d_i\bigr]\,\bigr\|_1 \;\ge\; \gamma ,
\]
where $\E_\lambda[d_i] = \sum_{\sminus} \lambda(\sminus)\, d_i(\sminus)$.
\end{theorem}

\begin{proof}
Feasibility for agent $i$ says
$\max_{\pay \in [0,1]^K} \min_{\sminus} \langle \pay, d_i(\sminus)\rangle \ge \gamma$.
The inner minimum is over a finite set, so it equals the minimum over its convex
hull. The map $(\pay,\lambda) \mapsto \langle \pay, \E_\lambda[d_i]\rangle$ is
bilinear and both $[0,1]^K$ and the simplex are convex and compact, so von
Neumann's minimax theorem gives
\[
  \max_{\pay} \min_{\lambda} \langle \pay, \E_\lambda[d_i]\rangle
  = \min_{\lambda} \max_{\pay} \langle \pay, \E_\lambda[d_i]\rangle
  = \min_{\lambda} \sum_y \bigl(\E_\lambda[d_i]\bigr)_y^{+},
\]
the last step because maximising a linear form over the unit cube sets
$\pay_y = 1$ on positive coordinates and $0$ elsewhere. Each $d_i(\sminus)$ is a
difference of probability vectors, so it sums to zero, as does any mixture; for a
zero-sum $v$, $\sum_y v_y^+ = \tfrac12\|v\|_1$.
\end{proof}

\paragraph{What is and is not new here.}
Both ingredients are classical: von Neumann's minimax theorem, and the
variational characterisation of total variation as the largest gap a bounded test
function can produce. The \emph{qualitative} content---that holding an agent to
account requires its deviation to move the distribution of the observed
signal---is also standard, and is exactly the role of the individual and pairwise
full-rank conditions of Fudenberg, Levine and Maskin~\cite{flm1994} in repeated
games with imperfect public monitoring. We claim neither. What
Theorem~\ref{thm:elim} adds is a \emph{quantitative, one-shot} form of that
condition---the shift must be at least $\gamma$ in total variation, robustly over
mixtures of teammate behaviour---and, decisively, the elimination of the payment
rule. Detectability conditions in the repeated-games literature are stated for a
\emph{given} monitoring structure; here the monitoring structure is the unknown,
and removing $\pay$ is what makes it possible to ask which structures a budget
can buy.

\begin{corollary}[Edge condition]
\label{cor:edge}
Taking $\lambda$ a point mass gives a necessary condition: for every agent $i$
and every teammate profile $\sminus$,
\[
  \TV\bigl(\enc[(1,\sminus),\cdot],\; \enc[(0,\sminus),\cdot]\bigr) \;\ge\; \gamma .
\]
\end{corollary}

\begin{corollary}[A bound mechanism design must respect]
\label{cor:effort}
Fix effort costs $c_i > 0$ and let $\gamma^\ast(\enc)$ denote the largest margin
the record $\enc$ implements. If $\gamma^\ast(\enc) < c_i$ then no bounded
payment rule computed from the record makes agent~$i$'s effort profitable
against every teammate profile: for \emph{every} $\pay_i \in [0,1]^K$ there is a
$\sminus$ with $\langle \pay_i, d_i(\sminus)\rangle < c_i$. Consequently no
mechanism whose transfers are measurable with respect to the record can sustain
agent~$i$'s effort, under any solution concept.
\end{corollary}

\begin{proof}
Immediate from Theorem~\ref{thm:elim}: $\gamma^\ast(\enc)$ is the value of the
minimax problem, so $\gamma^\ast(\enc) < c_i$ means no bounded $\pay_i$ achieves
$c_i$ uniformly in $\sminus$. The final clause holds because any such mechanism's
transfer to $i$ is a function of $Y$ alone, hence is some $\pay_i$ after
rescaling into $[0,1]$.
\end{proof}

The scope of Corollary~\ref{cor:effort} is worth stating plainly, since it is
what makes the rest of the paper a constraint on mechanism design rather than an
exercise in coding theory. It is a statement about \emph{information}, not about
equilibrium: it does not identify an equilibrium, does not model participation or
budget balance, and does not say what the designer should do. It says what no
designer can do. The quantities $\Dopt$, $\Kflat$ and the frontier computed below
are therefore feasibility boundaries---a record failing them is unusable for
settlement no matter how clever the mechanism placed on top of it.

\paragraph{Two further consequences.}
First, the problem becomes an optimisation over the record alone:
Theorem~\ref{thm:elim} replaces a bilinear search over $(\enc,\pay)$ by a
condition on $\enc$, which is what makes Section~\ref{sec:threshold} reachable---a
bilinear formulation did not terminate at $n = 7$, whereas the reduced question
is a feasibility problem over $2^n$ binary variables and runs in milliseconds.
Second, impossibility becomes cheap: if a record maps some profile and all $n$ of
its unilateral neighbours to the same label, every edge total variation there is
$0$, so by Corollary~\ref{cor:edge} no bounded payment achieves any positive
margin, for any decoder rather than merely a binary one.

Corollary~\ref{cor:edge} is the general statement of the phenomenon this paper is
about: \emph{a record supports accountability exactly to the extent that each
agent, acting alone, visibly moves it.} Reconstruction wants the opposite---it
merges profiles that are cheap to confuse in order to spend its $K$ labels well.
The rest of the paper measures the collision.
